\chapter{Conclusion}
\label{ch:conclusion}

This thesis establishes that an autonomous, multi-agent AI framework can successfully execute a rigorous, end-to-end validation of complex phenomenological physics models. By systematically evaluating $p_T$ spectra across ten multiplicity classes, we investigated the long-standing physical tension between soft thermal emission and hard scattering tails.

Our numerical pipeline exposed the intrinsic limitations of the proposed 2-component J\"uttner model. The model's failure to converge—driven by extreme parameter correlation and flat $\chi^2$ likelihood manifolds—demonstrates that na\"ively introducing additional moving thermal sources results in severe over-parameterization rather than physical insight. Conversely, our framework independently validated the Boltzmann-Gibbs Blast-Wave (BGBW) baseline, extracting freeze-out parameters that align perfectly with established ALICE experimental literature.

Ultimately, this work proves that resolving the soft/hard tension in high-energy collisions requires robust, mathematically stable formalisms (such as the Tsallis distribution or explicit soft/hard decompositions) rather than the proliferation of unconstrained thermal components. The AI-augmented methodology developed herein establishes a new paradigm for automated peer review, ensuring that phenomenological proposals meet rigorous standards of reproducibility and statistical stability in high-energy physics.
